Fix a feasible \((p,w)\) and a schedule in \(\mathcal C_{n,\beta}(B,\omega)\), and condition on that fixed schedule as required by Assumption~\(\mathrm{ass{:}fixed\text{-}finite\text{-}population}\).  Put \(\bar c_{\varnothing}=n^{-1}\sum_{i\in I_n}c_{i,\varnothing}\).  Assumption~\(\mathrm{ass{:}low\text{-}order\text{-}schedule}\) and independence in Assumption~\(\mathrm{ass{:}common\text{-}uniform\text{-}rollout}\) imply
\[
\mathbb E_U M(r)
=\bar c_{\varnothing}
+\sum_{S\in\mathcal S_{n,\beta}}a_Sr^{|S|},
\tag{13}
\]
because \(\mathbb E_U\prod_{j\in S}Z_j(r)=\Pr_U(U_j\le r\text{ for all }j\in S)=r^{|S|}\).  Definition~\(\mathrm{def{:}unbiased\text{-}weight\text{-}class}\) gives \(\sum_tw_t=0\) and \(\sum_tw_tp_t^d=1\) for every \(1\le d\le\beta\).  Applying these identities term by term in (13) yields
\[
\mathbb E_U\widehat\tau_{p,w}
=\sum_{S\in\mathcal S_{n,\beta}}a_S.
\tag{14}
\]
Separately, the low-order expansion gives the causal contrast
\[
\tau
=\frac1n\sum_{i\in I_n}\{Y_i(\mathbf1)-Y_i(\mathbf0)\}
=\sum_{S\in\mathcal S_{n,\beta}}a_S.
\tag{15}
\]
Equations (14)--(15) prove design unbiasedness; the equality with the causal estimand is derived from the potential outcomes rather than imposed by definition.

Lemma~\(\mathrm{lem{:}common\text{-}uniform\text{-}monomial\text{-}covariance}\) gives
\[
\operatorname{Var}_U(\widehat\tau_{p,w})=a^{\mathsf T}G(p,w)a.
\]
Since every \(\omega_d>0\), put \(b=\Omega^{-1/2}a\).  The ellipsoid assumption gives \(\lVert b\rVert_2^2\le B^2/n\), and the quadratic-form bound gives
\[
\operatorname{Var}_U(\widehat\tau_{p,w})
\le \frac{B^2}{n}
\lambda_{\max}\{\Omega^{1/2}G(p,w)\Omega^{1/2}\}
=V_{\mathrm{rob}}(p,w),
\tag{16}
\]
where the final equality is the fixed-design support-function conclusion of Theorem~\(\mathrm{thm{:}exact\text{-}covariance\text{-}saddle}\), together with Definition~\(\mathrm{def{:}variance\text{-}certificate}\).

Suppose first that \(V_{\mathrm{rob}}(p,w)>0\).  Set
\[
X=\widehat\tau_{p,w}-\tau,
\qquad
R=\sqrt{V_{\mathrm{rob}}(p,w)/\alpha}.
\]
Equations (14)--(16) give \(\mathbb E_UX=0\) and \(\mathbb E_UX^2\le V_{\mathrm{rob}}(p,w)\).  The pointwise inequality \(X^2\ge R^2\mathbf1\{|X|>R\}\) therefore implies
\[
\Pr_U(|X|>R)
\le\frac{\mathbb E_UX^2}{R^2}
\le\alpha.
\tag{17}
\]
Taking complements in (17) and applying Definition~\(\mathrm{def{:}chebyshev\text{-}interval}\) proves both displayed coverage formulations.  If \(V_{\mathrm{rob}}(p,w)=0\), then (16) and unbiasedness imply \(\mathbb E_UX^2=0\).  Since \(X^2\ge0\), this forces \(X=0\) almost surely, so the singleton interval covers with probability one.